% =====================================================================
%  Del Slop al Conocimiento Verificable
%  Universidad Icesi - Maestria en Inteligencia Artificial Aplicada
%  Formato: plantilla de presentacion de documentos, Facultad Barberi.
%  Compilar con: pdflatex informe.tex  (dos veces)
% =====================================================================
\documentclass[11pt,letterpaper]{article}

\usepackage[utf8]{inputenc}
\usepackage[T1]{fontenc}
\usepackage[spanish,es-tabla,es-nodecimaldot]{babel}
\usepackage{mathptmx}                 % Times New Roman
\usepackage[scaled=0.92]{helvet}
\usepackage{courier}

\usepackage[letterpaper,top=2.54cm,bottom=2.54cm,left=2.54cm,right=2.54cm,
            headheight=14pt,headsep=18pt]{geometry}

\usepackage{fancyhdr}
\usepackage{titlesec}
\usepackage{booktabs}
\usepackage{longtable}
\usepackage{array}
\usepackage{ragged2e}
\usepackage{caption}
\usepackage{tikz}
\usepackage{xcolor}
\usepackage[hidelinks]{hyperref}
\usepackage{xurl}

% ---------------------------------------------------------------- estilo
\definecolor{barfill}{HTML}{3B5C86}
\definecolor{barlight}{HTML}{9FB3CE}
\definecolor{hito}{HTML}{9B3A2E}
\definecolor{rejilla}{HTML}{C8CDD4}

\pagestyle{fancy}
\fancyhf{}
\fancyhead[L]{\footnotesize Universidad Icesi, Facultad Barberi de Ingeniería Diseño y Ciencias Aplicadas}
\fancyhead[R]{\footnotesize\thepage}
\renewcommand{\headrulewidth}{0.4pt}

\titleformat{\section}[block]{\centering\normalsize\scshape}{\Roman{section}.\quad}{0pt}{}
\titlespacing*{\section}{0pt}{16pt}{8pt}
\titleformat{\subsection}[block]{\normalsize\itshape}{\Alph{subsection}.\quad}{0pt}{}
\titlespacing*{\subsection}{\parindent}{11pt}{4pt}

\captionsetup[table]{labelsep=newline,justification=centering,font=sc,labelfont=sc,skip=2pt}
\captionsetup[figure]{labelsep=period,justification=raggedright,font=footnotesize,skip=6pt}
\renewcommand{\thetable}{\Roman{table}}
\renewcommand{\arraystretch}{1.18}

\setlength{\parindent}{1.3em}
\setlength{\parskip}{0pt}

\newcolumntype{L}[1]{>{\RaggedRight\arraybackslash}p{#1}}
\newcommand{\refhang}[1]{\par\noindent\hangindent=1.5em\hangafter=1 #1\par\vspace{2pt}}

% nota al pie sin marca, como en la plantilla
\newcommand{\notasinmarca}[1]{%
  \begingroup\renewcommand{\thefootnote}{}\footnote{#1}\addtocounter{footnote}{-1}\endgroup}

% ================================================================ inicio
\begin{document}

\thispagestyle{fancy}

\begin{center}
{\LARGE Del Slop al Conocimiento Verificable:\\[2pt]
Compilación incremental y trazabilidad de literatura científica\par}
\vspace{18pt}
{\ttfamily\small Juan Diego Ramírez\quad Juan David Cruz García\par}
{\ttfamily\small juan.ramirezzuluaga@u.icesi.edu.co, juan.cruz6@u.icesi.edu.co\par}
\end{center}

\vspace{20pt}

\notasinmarca{Documento recibido el 1 de septiembre de 2026. Proyecto de la Maestría en Inteligencia Artificial Aplicada, Universidad Icesi. Código y datos disponibles bajo licencia MIT en \url{https://github.com/1KVueltasAlCampo/llm-knowledge-management-wiki-agent}.}

\noindent\textbf{Resumen —} Las organizaciones pierden tiempo de trabajo buscando conocimiento
fragmentado y desactualizado, y los asistentes basados en Recuperación Aumentada por Generación
(RAG) que se adoptan para mitigarlo repiten el trabajo de síntesis en cada consulta, mezclan
documentos vigentes con obsoletos y no consolidan un activo reutilizable. Este trabajo propone un
pipeline MLOps que construye y mantiene una base de conocimiento en markdown interenlazado según el
patrón WikiLLM, trasladando el costo de síntesis desde la consulta hacia la ingesta. El problema se
instancia sobre literatura abierta de inteligencia artificial publicada en arXiv, donde la ausencia
de revisión por pares acumula texto generado por modelos de lenguaje sin evidencia verificable. La
evaluación ordena las fuentes por fecha, las ingiere de a una y guarda una copia versionada del wiki
tras cada ingesta; sobre cada copia se responde un conjunto de preguntas de control con respuesta
fijada de antemano, y se miden cobertura, fundamentación, resolución de contradicciones,
supervivencia de un error inyectado y costo, con una línea base RAG sobre el mismo corpus. La
contribución esperada es un protocolo de evaluación de mantenimiento incremental, su primer conjunto
de mediciones y una implementación abierta y reproducible. El documento presenta la formulación, la
metodología, el cronograma y los criterios de evaluación; los resultados experimentales corresponden
a una fase posterior.

\vspace{8pt}
\noindent\textbf{Índice de Términos —} Base de conocimiento, evaluación incremental, MLOps,
recuperación aumentada por generación.

% ====================================================================
\section{Introducción}

\subsection{Planteamiento del problema}

En las organizaciones, el conocimiento —procesos, decisiones técnicas, manuales, políticas— está
fragmentado y envejece rápido, y localizarlo consume tiempo de trabajo. Un estudio con 716 empleados
de cuatro entidades públicas encontró que el 22,3\,\% dedica cerca de media jornada semanal a buscar
la información que necesita, y que un 10,5\,\% dedica jornada y media, cerca del 30\,\% de su semana
(Nakash y Bouhnik, 2024). La cifra más citada en la industria, alrededor del 20\,\% de la semana
laboral, proviene de un informe anterior (Chui et al., 2012) y conviene leerla como orden de
magnitud. Ambas apuntan a una ineficiencia transversal y anterior a la adopción de inteligencia
artificial.

La respuesta habitual son asistentes basados en Recuperación Aumentada por Generación (RAG). La
evaluación sistemática de Han et al. (2025) muestra que ni el RAG vectorial ni los enfoques basados
en grafos dominan de forma absoluta: cada uno aventaja al otro según la naturaleza de la consulta.
Frente al problema de conocimiento organizacional, el RAG convencional presenta tres límites:

\begin{enumerate}
\item \textit{Costo operativo por consulta.} El modelo recupera, lee y sintetiza texto crudo cada vez
que alguien pregunta. El trabajo de síntesis se repite íntegro aunque la pregunta sea parecida a la
anterior (Karpathy, 2026).
\item \textit{Degradación de la coherencia.} El sistema recupera documentos obsoletos o en conflicto
con material reciente e intenta reconciliarlos en el momento de responder, lo que propaga errores y
produce la deriva que Myakala et al. (2026) formalizan como \textit{drift coherence}.
\item \textit{Incapacidad de consolidación.} El conocimiento sintetizado no persiste. La organización
no construye un activo que crezca, sino un buscador que vuelve a empezar de cero.
\end{enumerate}

\subsection{Dominio de estudio}

El proyecto instancia ese problema general sobre la literatura abierta de inteligencia artificial
publicada en arXiv. La elección es metodológica antes que temática: constituye la versión más
exigente del mismo problema y la única que puede evaluarse de forma reproducible y pública.

Desarrollar software con herramientas de inteligencia artificial exige consultar literatura abierta
de forma constante, y en este campo el intervalo entre publicar y construir se ha comprimido a
semanas. Estos repositorios, sin embargo, publican sin revisión por pares y acumulan \textit{slop}:
textos generados por modelos de lenguaje, correctamente redactados, sin código ni evidencia
verificable. El volumen hace inviable el filtrado manual —solo en febrero de 2026 arXiv recibió más
de 20.500 preprints de computación (arXiv, 2026a)— y las normas vigentes no cierran el vacío, porque
operan como sustitutos y no como filtros: la de octubre de 2025 restringe un formato de publicación
en lugar de detectar generación automática (arXiv, 2025), y la sanción de mayo de 2026 alcanza
únicamente el descuido evidente, como referencias inexistentes o instrucciones olvidadas dentro del
manuscrito (Dietterich, 2026). Elazar y Antoniak (2026) delimitan dónde queda el hueco: el texto
generado por modelos de lenguaje es casi seis veces más frecuente en artículos que no son revisiones,
justamente la categoría que ninguna norma alcanza.

RAG agrava este escenario en lugar de mitigarlo. Al recuperar por similitud semántica trata todos los
fragmentos como equivalentes y otorga la misma autoridad a un experimento validado que a una
afirmación sin sustento. En un corpus donde no puede presumirse que lo publicado fue verificado, esa
indiferencia deja de ser un detalle técnico: el juicio recae íntegramente en el lector, que no dispone
de metadato alguno para ejercerlo.

Adicionalmente, el corpus es de acceso libre, sin permisos ni costos de licencia, lo que permite
publicar datos y código y que un tercero replique el experimento. Un corpus corporativo real
impediría esa verificación.

\subsection{Pregunta de investigación y contribución esperada}

El trabajo se organiza en torno a una pregunta: \textit{¿cómo evoluciona la calidad de un WikiLLM a
medida que se le ingieren fuentes, y en qué condiciones supera a un sistema RAG tradicional?}

La formulación es deliberadamente abierta. Se busca establecer bajo qué condiciones compilar por
adelantado resulta conveniente, no demostrar que siempre lo es, dado que Han et al. (2025) ya
mostraron que ninguna arquitectura aventaja a las demás en toda tarea. El diseño incorpora además una
hipótesis de riesgo contra la propia propuesta: una base compilada puede blanquear el ruido,
convirtiendo una afirmación débil en prosa segura y permanente. Si eso ocurre, el WikiLLM resultaría
inferior a RAG en este escenario, y el trabajo lo reportará.

Los sistemas de referencia en generación de artículos enciclopédicos —STORM (Shao et al., 2024),
WikiChat (Semnani et al., 2023) y el conjunto de evaluación FreshWiki— miden artefactos generados una
sola vez. Ninguno evalúa lo que define al patrón WikiLLM: qué le ocurre a una base de conocimiento
mantenida en el tiempo. El patrón, en cambio, ya cuenta con implementaciones maduras; la más completa
es un complemento de Obsidian que ofrece detección de contradicciones, fusión de duplicados y páginas
protegidas contra sobrescritura (green-dalii, 2026). El problema de ingeniería, por tanto, está
resuelto, y la pregunta abierta se desplaza hacia la evaluación: la única cifra de desempeño publicada
en este espacio es un 27,1\,\% frente a un 24,1\,\% de la línea base, autorreportada por el propio
proyecto sobre su propio corpus, y midiendo recuperación en un wiki ya construido, no su conservación
durante el mantenimiento. Una ventaja de tres puntos sin protocolo independiente no establece que el
patrón funcione: establece que todavía no se ha medido lo relevante.

La contribución esperada consta de tres elementos: un protocolo de evaluación de mantenimiento
incremental, su primer conjunto de mediciones, y una implementación abierta que lo haga reproducible.

% ====================================================================
\section{Materiales y métodos}

\subsection{Marco teórico}

El patrón WikiLLM designa una base de conocimiento en texto plano —markdown interenlazado— que un
agente construye y mantiene de forma incremental a partir de fuentes crudas (Karpathy, 2026). Su
diferencia operativa con RAG reside en el momento en que ocurre la síntesis, con las consecuencias
que resume la Tabla~\ref{tab:comparacion}.

\begin{table}[htbp]
\centering
\caption{Comparación entre RAG y el patrón WikiLLM}
\label{tab:comparacion}
\begin{tabular}{L{3.6cm} L{4.9cm} L{4.9cm}}
\toprule
 & \textbf{RAG (línea base)} & \textbf{WikiLLM (propuesta)} \\
\midrule
Momento de la síntesis & En cada consulta & Una vez, al ingerir \\
Qué persiste & Fragmentos y vectores & Páginas redactadas y enlazadas \\
Manejo de la contradicción & Implícito, al responder & Explícito, marcado y fechado \\
Falla característica & Recuperar el fragmento equivocado & Propagar un error de síntesis anterior \\
\bottomrule
\end{tabular}
\end{table}

La arquitectura traslada el costo de síntesis desde la consulta hacia la ingesta: el agente compila
una vez y después responde sobre páginas ya redactadas, enlazadas y fechadas.

\subsection{Delimitación del corpus}

El corpus se restringe a una subárea de inteligencia artificial con tableros de clasificación
activos, entre 30 y 50 fuentes, en una ventana temporal posterior al corte de entrenamiento del
modelo evaluado. Esta delimitación cumple dos funciones. Las afirmaciones sobre resultados en
\textit{benchmarks} son numéricas, fechadas y objetivamente superables, lo que proporciona un criterio
de verdad sin juicio subjetivo. La ventana posterior al corte de entrenamiento impide que el modelo
responda de memoria en lugar de leer el wiki, condición sin la cual se estaría midiendo el
conocimiento previo del modelo y no la arquitectura.

La procedencia se trata como metadato de primera clase: cada afirmación conserva su fuente y sus
señales de respaldo —disponibilidad de código, presencia de experimentos propios, estado de
publicación y versión del preprint—, de modo que la calidad de la evidencia no se disuelva en la
prosa.

\subsection{Preprocesamiento}

La etapa de preparación de datos extrae texto de PDF y LaTeX conservando la estructura, normaliza
codificación y notación de cifras, elimina duplicados por huella de contenido y por versión del
preprint, unifica nombres de métodos y \textit{benchmarks}, y valida el formato descartando documentos
malformados antes de consumir una llamada al modelo.

\subsection{Elección del modelo}

La operación de ingesta debe leer una fuente nueva y revisar entre diez y quince páginas ya escritas.
Ese requisito funcional fija cuatro criterios de selección: ventana de contexto amplia, obediencia
estricta al formato de salida, costo por token predecible —el costo es una de las variables que el
experimento mide— y disponibilidad de una opción de ejecución local. La Tabla~\ref{tab:modelos}
recoge las alternativas consideradas.

\begin{table}[htbp]
\centering
\caption{Alternativas de modelo y función asignada}
\label{tab:modelos}
\begin{tabular}{L{3.9cm} L{1.5cm} L{1.5cm} L{1.5cm} L{4.3cm}}
\toprule
\textbf{Alternativa} & \textbf{Contexto} & \textbf{Formato} & \textbf{Local} & \textbf{Función} \\
\midrule
Comercial de gama alta & $\geq$200K & Alta & No & Ingesta y revisión \\
Comercial económico & $\geq$200K & Media-alta & No & Consultas y evaluación masiva \\
Abierto de ejecución local & 32--128K & Media & Sí & Control de contaminación \\
Vectores y reordenador & --- & --- & Sí & Únicamente la línea base RAG \\
\bottomrule
\end{tabular}
\end{table}

No se selecciona un modelo único, sino qué modelo interviene en cada etapa. El uso de \texttt{litellm}
permite que cambiar de proveedor sea una modificación de configuración y no de código, con lo que el
modelo se convierte en una variable del experimento y puede reportarse la relación entre costo y
calidad en lugar de un valor aislado.

\subsection{Arquitectura del pipeline}

El sistema se organiza como un flujo modular de seis pasos, cada uno reejecutable por separado sin
rehacer los anteriores: obtener la fuente, extraer afirmaciones con su origen, unificar entidades,
actualizar las páginas afectadas, revisar el wiki y guardar una copia versionada.

El modelo no opera como interfaz conversacional, sino como motor de un ciclo documental gobernado por
Git: lee una fuente, propone cambios sobre varias páginas simultáneamente y los entrega como
\textit{Pull Request}. Ninguna escritura llega al wiki sin aprobación humana previa. La diferencia con
las implementaciones existentes radica en la ubicación del control: hoy una página se protege
marcándola manualmente para que no sea sobrescrita, es decir, después de que el agente escribió. Aquí
la aprobación antecede a la escritura y queda registrada en el historial del repositorio, lo que
convierte la gobernanza en un dato medible —cuántos errores intercepta el revisor y a qué costo en
tiempo— en lugar de una garantía declarada.

El revisor evalúa los cambios en Obsidian, donde el grafo de enlaces muestra qué páginas toca cada
propuesta. La ingesta es idempotente: procesar dos veces la misma fuente no la duplica. El flujo
reintenta con espera creciente ante fallos del proveedor y guarda avance por fuente para reanudar sin
reprocesar el corpus completo.

\subsection{Herramientas}

Python con \texttt{litellm} para la abstracción de proveedores, \texttt{pydantic} para validar que la
salida del modelo tenga la forma esperada y \texttt{pytest} para detectar regresiones. La API de arXiv
provee fuentes y metadatos. Git y Obsidian conforman la capa de revisión humana sobre los mismos
archivos markdown, sin dependencia de plataforma. Las instrucciones del agente se versionan como
código y cada operación deja registro estructurado.

\subsection{Protocolo experimental}

Las fuentes se ordenan por fecha y se ingieren de a una. Tras cada ingesta se guarda una copia
versionada del wiki, y sobre cada copia se responde el mismo conjunto de sesenta preguntas de control,
cuya respuesta correcta fue fijada de antemano por los dos autores por separado, reportando el nivel
de acuerdo entre ambos. La comparación se establece contra una línea base RAG construida sobre el
mismo corpus.

Las ablaciones varían dos parámetros enunciados en la literatura y no medidos hasta ahora: la
frecuencia de la revisión automática del wiki y la profundidad de la intervención humana. Se practica
además inyección controlada de errores y de fuentes de baja calidad.

\subsection{Criterios de evaluación}

Se miden cinco variables sobre cada copia versionada, definidas en la Tabla~\ref{tab:metricas}. Las
métricas de resolución de contradicciones y supervivencia del error no son medidas estándar de
exactitud: evalúan si el sistema se corrige a sí mismo con el tiempo, que es el objetivo del producto.

\begin{table}[htbp]
\centering
\caption{Métricas de evaluación}
\label{tab:metricas}
\begin{tabular}{L{4.4cm} L{9.9cm}}
\toprule
\textbf{Métrica} & \textbf{Qué responde} \\
\midrule
Cobertura & ¿El wiki contiene la información que debería contener? \\
Fundamentación & ¿Cada afirmación se rastrea hasta su fuente, con su nivel de respaldo? \\
Resolución de contradicciones & ¿El wiki se corrige cuando un trabajo nuevo supera una cifra anterior? \\
Supervivencia del error & ¿Cuántas ingestas sobrevive un error introducido a propósito? \\
Costo & Tokens y tiempo por ingesta y por consulta, frente a la línea base RAG \\
\bottomrule
\end{tabular}
\end{table}

El conjunto de preguntas de control opera simultáneamente como prueba de regresión: si una
modificación del agente degrada la calidad, se detecta antes del despliegue. Dado que el sistema no
entrena pesos, el equivalente al reentrenamiento es la recompilación de las páginas que una fuente
nueva deja obsoletas.

% ====================================================================
\section{Resultados esperados}

Esta sección presenta el plan de ejecución y los criterios con los que se juzgarán los resultados. Las
mediciones corresponden a una fase posterior del proyecto y no se reportan aquí.

\subsection{Cronograma}

El semestre 2026-2 se ejecuta en siete fases sobre un equipo de dos personas. La
Figura~\ref{fig:gantt} representa su distribución temporal y los dos puntos de corte que gobiernan la
contingencia. El cronograma es una propuesta de trabajo y no un compromiso cerrado: las fechas se
ajustarán tras la delimitación del alcance con el tutor y según el calendario académico oficial.

\begin{figure}[htbp]
\centering
\begin{tikzpicture}[x=0.135cm, y=-0.60cm, font=\footnotesize]

% --- rejilla semanal
\foreach \d in {0,7,14,21,28,35,42,49,56,63,70,77,84}{
  \draw[rejilla, line width=0.3pt] (\d,-0.45) -- (\d,7.35);
}
% --- limites de mes
\foreach \d in {0,30,61,88}{
  \draw[gray!70, line width=0.6pt] (\d,-0.75) -- (\d,7.35);
}
% --- etiquetas de mes
\node[anchor=south] at (15,-0.80) {\textsc{septiembre}};
\node[anchor=south] at (45.5,-0.80) {\textsc{octubre}};
\node[anchor=south] at (74.5,-0.80) {\textsc{noviembre}};

% --- eje
\draw[gray!70, line width=0.6pt] (0,-0.45) -- (88,-0.45);

% --- barras
\newcommand{\fase}[4]{%
  \node[anchor=east, font=\footnotesize] at (-2,#1) {#4};
  \fill[barfill] (#2,#1-0.19) rectangle (#3,#1+0.19);
}
\fase{0.6}{0}{7}{\textbf{F0} Delimitación}
\fase{1.6}{7}{21}{\textbf{F1} Corpus}
\fase{2.6}{21}{35}{\textbf{F2} Protocolo}
\fase{3.6}{35}{49}{\textbf{F3} Agente WikiLLM v1}
\fase{4.6}{49}{63}{\textbf{F4} Línea base}
\fase{5.6}{63}{77}{\textbf{F5} Experimento}
\fase{6.6}{77}{88}{\textbf{F6} Cierre}

% --- hitos de control
\fill[hito] (35,2.6) circle (2.1pt);
\fill[hito] (63,4.6) circle (2.1pt);
\draw[hito, dashed, line width=0.5pt] (35,2.6) -- (35,7.15);
\draw[hito, dashed, line width=0.5pt] (63,4.6) -- (63,7.15);
\node[hito, anchor=north west, font=\scriptsize] at (35.5,7.15) {5 oct};
\node[hito, anchor=north west, font=\scriptsize] at (63.5,7.15) {2 nov};

\end{tikzpicture}
\caption{Cronograma de ejecución del semestre 2026-2. Los marcadores en rojo señalan los dos puntos de
corte: cierre de F2 el 5 de octubre y cierre de F4 el 2 de noviembre.}
\label{fig:gantt}
\end{figure}

\begin{table}[htbp]
\centering
\caption{Entregables verificables por fase}
\label{tab:fases}
\begin{tabular}{L{3.0cm} L{2.0cm} L{7.2cm} L{2.2cm}}
\toprule
\textbf{Fase} & \textbf{Fechas} & \textbf{Entregable verificable} & \textbf{Responsable} \\
\midrule
F0 Delimitación y aprobación & 1--7 sep & Acta de alcance validada con el tutor: subárea, ventana temporal y criterio de inclusión del corpus. Repositorio público con licencia MIT. & Ambos \\
F1 Corpus e ingeniería de datos & 8--21 sep & Corpus versionado de 30 a 50 preprints con extracción de texto y estructura. Base de metadatos y señales de calidad por fuente. Reporte de calidad de datos. & J. D. Ramírez \\
F2 Protocolo de evaluación & 22 sep--5 oct & Sesenta preguntas de control con respuesta verificable en las fuentes, respondidas por separado por ambos autores y con el nivel de acuerdo reportado. Script de evaluación automática. & Ambos \\
F3 Agente WikiLLM v1 & 6--19 oct & Operaciones de ingesta y revisión funcionando. Wiki compilado sobre el corpus con procedencia por afirmación. Flujo de revisión humana vía \textit{Pull Request}. & J. D. Cruz \\
F4 Línea base y orquestación & 20 oct--2 nov & Línea base RAG reproducible sobre el mismo corpus. Pipeline orquestado con ejecución programada, reintentos, registro estructurado y monitoreo de costo. & Ambos \\
F5 Experimento & 3--16 nov & Ingestas incrementales con copia versionada por ingesta, tres ablaciones de revisión, inyección controlada de errores y de fuentes de baja calidad. Resultados. & Ambos \\
F6 Cierre & 17--27 nov & Informe final, documentación de reproducción, artículo comparativo y sustentación oral. & Ambos \\
\bottomrule
\end{tabular}
\end{table}

\subsection{Hitos de control y contingencia}

El riesgo dominante del plan es que F3 se desplace y comprima F5, que es donde reside la
contribución. Se fijan dos puntos de corte. El 5 de octubre debe estar cerrada F2; si las preguntas de
control no están listas, se reducen de sesenta a treinta antes que retrasar el experimento, dado que
sin preguntas de control no hay nada que medir: la fase no es negociable en existencia, solo en
tamaño. El 2 de noviembre debe estar cerrada F4; si el pipeline orquestado no está estable, el
experimento se ejecuta con disparo manual y la orquestación se documenta como limitación. En ambos
casos se protege el experimento y no la infraestructura.

El alcance recortable sigue un orden establecido de antemano: primero las ablaciones de intervención
humana, después las de frecuencia de revisión automática y por último el número de fuentes. La línea
base RAG y la inyección controlada de errores no son recortables, porque sin ellas el resultado no
resulta interpretable.

\subsection{Criterios de éxito}

El trabajo se considera exitoso si produce, sobre el protocolo descrito, una curva de evolución de las
cinco métricas a lo largo de las ingestas sucesivas, con la comparación contra la línea base RAG y con
el punto de equilibrio de costo expresado en número de consultas por fuente ingerida. Un resultado
desfavorable al WikiLLM constituye igualmente un resultado válido, en tanto delimita las condiciones
en que cada arquitectura conviene.

% ====================================================================
\section{Discusión}

\subsection{Riesgos del diseño}

El riesgo central es la hipótesis de blanqueo enunciada en la Sección I: que la compilación convierta
literatura sin evidencia en prosa autoritativa. En este dominio el riesgo es mayor que en un corpus
revisado por pares, precisamente porque la relación entre señal y ruido de la fuente de entrada es
decreciente. No se declara como salvedad: se mide, inyectando fuentes de baja calidad y observando con
cuánta seguridad el wiki termina afirmándolas.

Un segundo riesgo es la propagación de errores de síntesis. En RAG el daño de una recuperación
incorrecta dura una respuesta; en una base compilada queda escrito de forma permanente y puede
sobrevivir a múltiples ingestas posteriores. La métrica de supervivencia del error existe para
cuantificarlo.

\subsection{Riesgos éticos y su mitigación}

\textit{Seguridad.} Un documento fuente puede contener instrucciones dirigidas al agente. Como se
señaló, en RAG el efecto se limita a una respuesta, mientras que aquí quedaría inscrito en el wiki de
forma persistente. La mitigación combina tres medidas: las fuentes se tratan siempre como datos y
nunca como órdenes, toda escritura pasa por revisión humana previa, y el sistema se prueba con
documentos preparados para inducir ese comportamiento, reportando la proporción que se detiene.

\textit{Sesgo.} Se atienden tres frentes: el criterio de inclusión de fuentes, declarado
explícitamente; la tendencia a que lo más reciente sobrescriba lo anterior sin evidencia suficiente; y
el blanqueo descrito en el apartado previo.

\textit{Privacidad.} El corpus es público y carece de datos personales. Todo texto enviado a un
proveedor externo queda registrado, y la variante con modelo local permite ejecutar el pipeline
completo sin que los datos abandonen el equipo, midiendo cuánta calidad cuesta esa opción.

La revisión humana atraviesa los tres frentes y su eficacia se reporta como dato: errores detectados
sobre errores introducidos, y minutos de revisión por error detectado.

\subsection{Limitaciones previstas}

El alcance es de un semestre y un solo dominio, con un corpus de entre 30 y 50 fuentes. La
contribución esperada es el protocolo y su primer conjunto de mediciones, no una respuesta definitiva
sobre la conveniencia del patrón. La generalización a corpus corporativos, con permisos y datos
sensibles, queda fuera del alcance y constituye una línea de continuación.

El conjunto de preguntas de control es construido por los propios autores, lo que introduce un riesgo
de sesgo de diseño que se mitiga parcialmente respondiéndolas por separado y reportando el nivel de
acuerdo, pero no se elimina. Un protocolo maduro requeriría anotadores independientes.

% ====================================================================
\section{Conclusiones}

El conocimiento organizacional fragmentado impone un costo medible, y los asistentes RAG que se
adoptan para reducirlo repiten el trabajo de síntesis, mezclan material vigente con obsoleto y no
consolidan un activo. El patrón WikiLLM propone trasladar la síntesis de la consulta a la ingesta,
pero su propiedad definitoria —el mantenimiento de la base en el tiempo— no ha sido evaluada por la
literatura disponible, que mide artefactos generados una sola vez.

Este trabajo delimita ese vacío y propone un protocolo para cubrirlo: ingesta incremental con copia
versionada, un conjunto fijo de preguntas de control con respuesta establecida de antemano, cinco
métricas que incluyen la corrección del sistema sobre sí mismo, y una línea base RAG sobre el mismo
corpus. El diseño incorpora una hipótesis de riesgo contra la propia propuesta y se compromete a
reportarla si se confirma.

Los resultados experimentales corresponden a la fase siguiente. La aplicación futura más inmediata es
la extensión del protocolo a corpus organizacionales, donde el problema original se manifiesta con
restricciones de permisos y privacidad que el corpus público no presenta.

% ====================================================================
\section*{Apéndice: cumplimiento de la rúbrica}

\small
\begin{longtable}{L{2.9cm} L{11.4cm}}
\toprule
\textbf{Criterio} & \textbf{Qué hace el proyecto} \\
\midrule
\endfirsthead
\toprule
\textbf{Criterio} & \textbf{Qué hace el proyecto} \\
\midrule
\endhead
\bottomrule
\endfoot
Cr1 Formulación & Se plantea como problema que las organizaciones pierden hasta un 30\,\% de la jornada buscando información y que los asistentes RAG actuales no lo resuelven, porque repiten el trabajo en cada consulta, mezclan lo vigente con lo obsoleto y no acumulan nada. El caso se estudia sobre literatura científica de IA, donde además ya no se distingue lo verificado del \textit{slop}. \\
Cr2 Elección del modelo & Se comparan cuatro tipos de modelo —comercial de gama alta, comercial económico, abierto local y un buscador por vectores para la línea base— frente a cuatro requisitos: ventana de contexto amplia, obediencia al formato de salida, costo por token y posibilidad de ejecución local. \\
Cr3 Integración & Se integra el modelo como motor de un flujo documental sobre Git: lee una fuente, propone cambios en varias páginas y los entrega como \textit{Pull Request}. Ningún cambio entra al wiki sin aprobación humana. \\
Cr4 Evaluación & Cinco métricas medidas después de cada ingesta contra un conjunto de preguntas cuya respuesta correcta se fijó de antemano: cobertura, fundamentación, resolución de contradicciones, supervivencia del error y costo. \\
Cr5 Tareas automatizables & Se automatiza leer, resumir y cruzar fuentes, que es la parte repetitiva y costosa en tiempo. Se deja al humano únicamente el juicio sobre qué información es válida. \\
Cr6 Eficiencia & Se mide el costo del sistema frente a un RAG tradicional sobre el mismo corpus, y se calcula a partir de cuántas consultas por fuente la compilación previa resulta más barata. \\
Cr7 Riesgos éticos & Seguridad: una fuente puede traer instrucciones ocultas que quedarían escritas de forma permanente, de modo que se prueba el sistema con documentos preparados para engañarlo y se reporta cuántos se detienen. Sesgo: se mide si el sistema convierte literatura sin evidencia en texto que suena confiable. Privacidad: el corpus es público y todo puede ejecutarse con un modelo local. \\
Cr8 Documentación & Se publica el repositorio bajo licencia MIT, con guía de reproducción paso a paso, las instrucciones del agente versionadas y un registro de las decisiones de diseño. \\
R.7.1 Cr1 Preprocesamiento & Se construye una etapa que extrae texto de PDF y LaTeX conservando la estructura, normaliza formatos, elimina duplicados y versiones repetidas del mismo preprint, y descarta documentos malformados antes de gastar una llamada al modelo. \\
R.7.1 Cr2 Herramientas & Python con \texttt{litellm} para cambiar de proveedor sin tocar código, \texttt{pydantic} para verificar que la salida del modelo tenga la forma esperada y \texttt{pytest}, más la API de arXiv y Git junto a Obsidian para la revisión humana. \\
R.7.1 Cr3 Pipeline & Un pipeline de seis pasos: traer la fuente, extraer afirmaciones con su origen, unificar entidades, actualizar páginas, revisar el wiki y guardar una copia versionada, cada uno reejecutable por separado. \\
R.7.2 Cr4 Despliegue & Se automatiza la ejecución del pipeline por calendario o a demanda, con seguimiento del costo y el tiempo de cada paso y aviso automático si la calidad se degrada. \\
R.7.2 Cr5 Robustez & La ingesta es idempotente, con reintentos ante fallos del proveedor y avance guardado por fuente. Se reporta cómo crece el tiempo de respuesta a medida que crece el wiki. \\
R.7.2 Cr6 MLOps & Se versionan las instrucciones del agente como código y se registra cada operación; las preguntas de evaluación funcionan como prueba de regresión. Como el sistema no entrena pesos, el equivalente al reentrenamiento es recompilar las páginas que una fuente nueva deja obsoletas. \\
RA1 y RA2 & El documento se estructura en las secciones solicitadas, con tablas donde la comparación es el contenido y citas APA verificadas en la fuente original. La sustentación se organiza sobre una sola tesis, con respuesta preparada para las objeciones previsibles. \\
\end{longtable}
\normalsize

% ====================================================================
\section*{Referencias}

\small

\refhang{arXiv. (2025, 31 de octubre). \textit{Attention authors: Updated practice for review articles and position papers in arXiv CS category}. arXiv Blog. \url{https://blog.arxiv.org/2025/10/31/attention-authors-updated-practice-for-review-articles-and-position-papers-in-arxiv-cs-category/}}

\refhang{arXiv. (2026a). \textit{Monthly submissions}. \url{https://arxiv.org/stats/monthly_submissions}}

\refhang{Chui, M., Manyika, J., Bughin, J., Dobbs, R., Roxburgh, C., Sarrazin, H., Sands, G., y Westergren, M. (2012). \textit{The social economy: Unlocking value and productivity through social technologies}. McKinsey Global Institute. \url{https://www.mckinsey.com/industries/technology-media-and-telecommunications/our-insights/the-social-economy}}

\refhang{Dietterich, T. (2026, 15 de mayo). \textit{Updated practice for papers containing unverified LLM output in arXiv CS category}. arXiv Blog.}

\refhang{Elazar, Y., y Antoniak, M. (2026). \textit{LLM-generated or human-written? Comparing review and non-review papers on arXiv} (arXiv:2601.17036). arXiv. \url{https://arxiv.org/abs/2601.17036}}

\refhang{green-dalii. (2026). \textit{Karpathy LLM Wiki: Obsidian plugin} [Software]. GitHub. \url{https://github.com/green-dalii/obsidian-llm-wiki}}

\refhang{Han, H., Ma, L., Wang, Y., Shomer, H., Lei, Y., Qi, Z., Guo, K., Hua, Z., Long, B., Liu, H., Aggarwal, C. C., y Tang, J. (2025). \textit{RAG vs. GraphRAG: A systematic evaluation and key insights} (arXiv:2502.11371). arXiv. \url{https://arxiv.org/abs/2502.11371}}

\refhang{Karpathy, A. (2026, 4 de abril). \textit{LLM wiki: A pattern for building personal knowledge bases using LLMs} [Gist]. GitHub. \url{https://gist.github.com/karpathy/442a6bf555914893e9891c11519de94f}}

\refhang{Myakala, P. K., Agrawal, M., y Manche, R. (2026). \textit{BeliefShift: Benchmarking temporal belief consistency and opinion drift in LLM agents} (arXiv:2603.23848). arXiv. \url{https://arxiv.org/abs/2603.23848}}

\refhang{Nakash, M., y Bouhnik, D. (2024). How much time does the workforce spend searching for information in the ``new normal''? En \textit{iConference 2024 Proceedings}. \url{https://www.ideals.illinois.edu/items/129980}}

\refhang{Semnani, S. J., Yao, V. Z., Zhang, H. C., y Lam, M. S. (2023). WikiChat: Stopping the hallucination of large language model chatbots by few-shot grounding on Wikipedia. En \textit{Findings of the Association for Computational Linguistics: EMNLP 2023} (pp. 2387--2413). \url{https://arxiv.org/abs/2305.14292}}

\refhang{Shao, Y., Jiang, Y., Kanell, T. A., Xu, P., Khattab, O., y Lam, M. S. (2024). Assisting in writing Wikipedia-like articles from scratch with large language models. En \textit{Proceedings of NAACL 2024}. \url{https://arxiv.org/abs/2402.14207}}

\normalsize

\end{document}
